\chapter{Reference Elements, Basis Functions, and Quadrature}
\label{ch:refelement}

\chapterorient{The previous chapter named four function spaces and told you which
field belongs in each, but it left the basis functions themselves as abstractions:
``the hat function on a vertex,'' ``the edge function with unit circulation.'' This
chapter opens the engine room. We will see how a basis function is actually
\emph{represented} inside a finite-element code, how it is \emph{integrated} to fill
in a matrix entry, and why both jobs are done not on the physical mesh but on a
single canonical \emph{reference element}. Three ideas carry the chapter: shape
functions defined once on a reference cell, a geometric map that carries them to
each physical cell, and a quadrature rule that turns the integral defining each
matrix entry into a finite weighted sum. By the end you will be able to write the
hat functions of a triangle in closed form, transform a derivative through a
Jacobian, integrate a product of basis functions by Gauss quadrature and check it
against the exact answer, and read the element-local tensor whose axes
($E,L,P,C,G$) the matrix-free pipeline of \cref{ch:matrixfree} runs on. This is the
layer where the finite-element method stops being a theory and becomes arithmetic.}

\section{One element, tabulated once}

A mesh of an engineering device has thousands to millions of cells. Each cell needs
its own basis functions, its own integrals, its own contribution to the global
matrix. If we defined a fresh set of polynomials on every physical cell---fitting
them to that cell's particular vertices, recomputing their integrals from
scratch---the cost would be enormous and the bookkeeping hopeless. Worse, it would
be \emph{wasteful}, because the cells are not really different: a tetrahedron is a
tetrahedron, whatever its size or orientation. The shape of a linear basis function
on one tetrahedron is, in a precise sense we are about to make, the \emph{same} as on
any other.

The finite-element method exploits this with a single organizing decision: define
the basis functions \emph{once}, on a fixed canonical cell called the
\emph{reference element}, and reach every physical cell by a geometric map. The
reference element $\hat K$ is a cell of standard shape in its own coordinate
system---we use the \emph{reference triangle} $\hat K = \{(\hat x,\hat y): \hat
x\ge 0,\ \hat y\ge 0,\ \hat x+\hat y\le 1\}$ in two dimensions and the
\emph{reference tetrahedron} in three. On $\hat K$ we lay down a family of
\emph{shape functions} $\hat\phi_1,\dots,\hat\phi_L$ once and for all, and we
\emph{tabulate} them: we compute and store their values, and the values of their
derivatives, at a fixed set of points. That single table serves the entire mesh.

\begin{keyidea}
The reference element is the finite-element method's central act of reuse. Shape
functions are defined and tabulated \emph{once}, on a canonical reference cell
$\hat K$; each physical cell $K$ is reached by a geometric map $\vx = F_K(\hat\vx)$.
A basis function on $K$ is just a reference shape function composed with the inverse
map. One tabulation---a small table of basis values and derivatives at a few
points---drives the integration of every cell in the mesh. The mesh contributes only
its geometry (through the map $F_K$); the polynomial content is shared.
\end{keyidea}

\Cref{fig:refmap} shows the arrangement. The reference cell sits on the left in its
own $(\hat x,\hat y)$ coordinates; the map $F_K$ carries it to a physical cell $K$
somewhere in the mesh. Everything hard---the choice of polynomials, their values,
their derivatives---lives on the left, computed once. Everything cell-specific---the
size, the shape, the orientation---lives in the map $F_K$, which is cheap to apply.
This split is not merely an optimization; it is the conceptual backbone of the whole
construction, and it recurs at every level. When \cref{ch:matrixfree} factors the
operator into $\mathsf{G}$ (gather to elements), $\mathsf{B}$ (apply the tabulated
basis), $\mathsf{D}$ (apply the per-point geometry and material), and their
transposes, the tabulated basis $\mathsf{B}$ is exactly the reference shape
functions of this chapter, and the geometry factor $\mathsf{D}$ is exactly the map's
Jacobian.

\begin{figure}[t]
\centering
\begin{tikzpicture}[scale=1.0, every node/.style={font=\scriptsize}]
% --- reference triangle on the left ---
\begin{scope}
  \node[font=\footnotesize\bfseries] at (1.0,2.35) {reference $\hat K$};
  \fill[simBgBlue] (0,0)--(2,0)--(0,2)--cycle;
  \draw[simInk,thick] (0,0)--(2,0)--(0,2)--cycle;
  \fill[simBlue] (0,0) circle (2.2pt) node[below left]{$\hat\vx_0$};
  \fill[simBlue] (2,0) circle (2.2pt) node[below right]{$\hat\vx_1$};
  \fill[simBlue] (0,2) circle (2.2pt) node[above left]{$\hat\vx_2$};
  \draw[->,simGray] (-0.35,-0.35)--(0.7,-0.35) node[right]{$\hat x$};
  \draw[->,simGray] (-0.35,-0.35)--(-0.35,0.7) node[above]{$\hat y$};
  \node[simGray] at (0.7,0.6) {$\hat\phi_i$ tabulated here};
\end{scope}
% --- the map ---
\draw[flow, simRust, thick] (2.6,1.0) .. controls (3.4,1.5) and (4.0,1.5) .. (4.8,1.0)
  node[midway, above, font=\scriptsize]{$\vx = F_K(\hat\vx)$};
% --- physical cell on the right (skewed) ---
\begin{scope}[shift={(5.4,0)}]
  \node[font=\footnotesize\bfseries] at (1.3,2.35) {physical $K$};
  \fill[simBgRust] (0.2,0.1)--(2.4,0.5)--(1.0,2.0)--cycle;
  \draw[simInk,thick] (0.2,0.1)--(2.4,0.5)--(1.0,2.0)--cycle;
  \fill[simRust] (0.2,0.1) circle (2.2pt) node[below left]{$\vx_0$};
  \fill[simRust] (2.4,0.5) circle (2.2pt) node[below right]{$\vx_1$};
  \fill[simRust] (1.0,2.0) circle (2.2pt) node[above]{$\vx_2$};
  \draw[->,simGray] (-0.2,-0.35)--(0.85,-0.35) node[right]{$x$};
  \draw[->,simGray] (-0.2,-0.35)--(-0.2,0.7) node[above]{$y$};
  \node[simGray] at (1.2,0.7) {$\phi_i = \hat\phi_i\circ F_K^{-1}$};
\end{scope}
\end{tikzpicture}
\caption[The reference element and the geometric map]{The organizing picture of the
chapter. Shape functions $\hat\phi_i$ are defined and tabulated once on the reference
cell $\hat K$ (left, in its own $(\hat x,\hat y)$ coordinates). The geometric map
$F_K$ carries $\hat K$ onto each physical cell $K$ of the mesh (right). A physical
basis function is the reference shape function pulled back through the map,
$\phi_i = \hat\phi_i\circ F_K^{-1}$. The mesh supplies only the map; the polynomials
are shared across every cell. This is the $\hat\phi\to\mathsf{B}$, $F_K\to\mathsf{D}$
split that \cref{ch:matrixfree} runs on.}
\label{fig:refmap}
\end{figure}

The rest of the chapter fills in the three pieces. \Cref{sec:shape} builds the shape
functions on $\hat K$. \Cref{sec:map} builds the map $F_K$ and shows how it
transforms values, derivatives, and volumes. \Cref{sec:quad} replaces the integrals
by quadrature. \Cref{sec:vector} notes how the vector spaces of
\cref{ch:functionspaces} fit the same frame. \Cref{sec:tensor} collects the data
into the element-local tensor. Two worked examples (\cref{wex:massmatrix},
\cref{wex:jacobian}) carry the arithmetic all the way through.

\section{Nodal shape functions on the reference simplex}
\label{sec:shape}

We begin with the simplest and most important family: the \emph{nodal} or
\emph{Lagrange} shape functions, the basis of the $\Hone$ space of
\cref{ch:functionspaces}. These are the building blocks of every scalar potential
problem and the scaffolding from which the vector elements are built.

\subsection{Barycentric coordinates}

The natural coordinate system on a simplex (a triangle, a tetrahedron) is not
Cartesian but \emph{barycentric}. On the reference triangle with vertices
$\hat\vx_0,\hat\vx_1,\hat\vx_2$, the barycentric coordinates
$\lambda_0,\lambda_1,\lambda_2$ of a point $\hat\vx$ are the unique weights, summing
to one, that express $\hat\vx$ as a weighted average of the vertices:
\[
  \hat\vx = \lambda_0\hat\vx_0 + \lambda_1\hat\vx_1 + \lambda_2\hat\vx_2,
  \qquad \lambda_0+\lambda_1+\lambda_2 = 1.
\]
Geometrically, $\lambda_i$ is the fractional area of the sub-triangle opposite
vertex $i$: it is $1$ at vertex $i$, $0$ on the opposite edge, and falls off linearly
in between. For the standard reference triangle with $\hat\vx_0=(0,0)$,
$\hat\vx_1=(1,0)$, $\hat\vx_2=(0,1)$ they have the closed form
\[
  \lambda_0 = 1-\hat x-\hat y, \qquad
  \lambda_1 = \hat x, \qquad
  \lambda_2 = \hat y.
\]
The barycentric coordinates \emph{are} the lowest-order nodal shape functions: each
$\lambda_i$ is a linear polynomial equal to $1$ at its own vertex and $0$ at the
other two. We write $\hat\phi_i = \lambda_i$ for the $P_1$ (degree-one) basis. The
defining property is the \emph{cardinality} (or \emph{interpolation}, or
\emph{unisolvence}) relation from \cref{ch:functionspaces},
\begin{equation}
  \hat\phi_i(\hat\vx_j) = \delta_{ij},
  \label{eq:cardinality}
\end{equation}
``the $i$-th shape function is one at its own node and zero at every other node.''
This single relation is what makes the coefficients of a finite-element field equal
to its \emph{nodal values}: if $u_h = \sum_j c_j\hat\phi_j$, then evaluating at node
$i$ and using \eqref{eq:cardinality} gives $u_h(\hat\vx_i)=c_i$. The coefficient on
$\hat\phi_i$ is literally the field's value at node $i$---the DoF functional
$\sigma_i(u)=u(\hat\vx_i)$ of \cref{ch:functionspaces} reads it directly.

\begin{notationbox}
We mark reference-cell quantities with a hat: $\hat\vx$ for a reference point,
$\hat\phi_i$ for a reference shape function, $\hat K$ for the reference cell, and
$\Grad_{\hat\vx}$ for a gradient in reference coordinates. Unhatted symbols
($\vx,\phi_i,K,\Grad_\vx$) live on the physical cell. The map $F_K$ and its Jacobian
$J$ are the bridge between the two, and keeping the hats straight is the single most
common source of sign and transpose errors in finite-element code.
\end{notationbox}

\Cref{fig:reftet} shows the reference triangle and tetrahedron with their node
numbering. The triangle has three vertices (three $P_1$ nodes); the tetrahedron has
four. On the tetrahedron the barycentric coordinates are
$\lambda_0=1-\hat x-\hat y-\hat z$, $\lambda_1=\hat x$, $\lambda_2=\hat y$,
$\lambda_3=\hat z$---the same pattern, one more coordinate.

\begin{figure}[t]
\centering
\begin{tikzpicture}[scale=1.0, every node/.style={font=\scriptsize}]
% --- reference triangle with barycentric labels ---
\begin{scope}
  \node[font=\footnotesize\bfseries] at (1.1,2.55) {reference triangle};
  \fill[simBgBlue] (0,0)--(2.2,0)--(0,2.2)--cycle;
  \draw[simInk,thick] (0,0)--(2.2,0)--(0,2.2)--cycle;
  \fill[simBlue] (0,0)   circle (2.4pt) node[below left] {$0$};
  \fill[simBlue] (2.2,0) circle (2.4pt) node[below right]{$1$};
  \fill[simBlue] (0,2.2) circle (2.4pt) node[above left] {$2$};
  \node[simGray] at (0.35,0.35) {$\lambda_0$};
  \node[simGray] at (1.55,0.28) {$\lambda_1$};
  \node[simGray] at (0.3,1.55)  {$\lambda_2$};
  \node[simInk] at (1.1,-0.55) {$\lambda_0=1{-}\hat x{-}\hat y,\ \lambda_1=\hat x,\ \lambda_2=\hat y$};
\end{scope}
% --- reference tetrahedron ---
\begin{scope}[shift={(5.2,0)}]
  \node[font=\footnotesize\bfseries] at (1.1,2.55) {reference tetrahedron};
  \coordinate (n0) at (0,0);
  \coordinate (n1) at (2.2,0.3);
  \coordinate (n2) at (0.9,1.9);
  \coordinate (n3) at (1.0,0.7);
  \fill[simBgBlue,opacity=0.6] (n0)--(n1)--(n2)--cycle;
  \draw[simGray,densely dashed] (n0)--(n3) (n1)--(n3) (n2)--(n3);
  \draw[simInk,thick] (n0)--(n1)--(n2)--cycle;
  \fill[simBlue] (n0) circle (2.4pt) node[below left] {$0$};
  \fill[simBlue] (n1) circle (2.4pt) node[below right]{$1$};
  \fill[simBlue] (n2) circle (2.4pt) node[above]      {$2$};
  \fill[simBlue] (n3) circle (2.4pt) node[right]      {$3$};
  \node[simInk] at (1.1,-0.55) {$4$ vertices, $6$ edges, $4$ faces};
\end{scope}
\end{tikzpicture}
\caption[Reference simplices with node numbering]{The reference simplices and their
barycentric coordinates. \emph{Left:} the reference triangle, vertices numbered
$0,1,2$; the barycentric coordinate $\lambda_i$ is the $P_1$ shape function that
equals $1$ at node $i$ and falls linearly to $0$ on the opposite edge. \emph{Right:}
the reference tetrahedron, vertices $0$--$3$ (back edges dashed), with four
barycentric coordinates following the same pattern. These node counts---$3$ for the
triangle, $4$ for the tetrahedron---are the lowest-order $\Hone$ DoF counts of
\cref{ch:functionspaces}.}
\label{fig:reftet}
\end{figure}

\subsection{Higher order: the $P_p$ families}

The $P_1$ basis represents a field as piecewise-linear: exact on linear fields,
first-order accurate on smooth ones. To do better we raise the polynomial degree.
The degree-$p$ Lagrange space $P_p$ on the triangle consists of all bivariate
polynomials of total degree $\le p$; it has
\[
  \dim P_p = \binom{p+2}{2} = \frac{(p+1)(p+2)}{2}
\]
shape functions (the count of monomials $\hat x^a\hat y^b$ with $a+b\le p$). For
$p=1$ that is $3$ (the vertices); for $p=2$ it is $6$. The $P_2$ basis adds one node
at the midpoint of each edge, so its six nodes are the three vertices plus the three
edge midpoints. Each $P_2$ shape function is still defined by the cardinality
relation \eqref{eq:cardinality}---one at its own node, zero at the other five---and
each has a clean barycentric form. The three vertex functions are
$\hat\phi_i = \lambda_i(2\lambda_i-1)$, and the three edge-midpoint functions are
$\hat\phi_{ij} = 4\lambda_i\lambda_j$. (You can verify directly that
$\lambda_0(2\lambda_0-1)$ is $1$ at vertex $0$, where $\lambda_0=1$, and $0$ at every
other node, where either $\lambda_0=0$ or $\lambda_0=\tfrac12$.)

\Cref{fig:shapefns} plots the $P_1$ and $P_2$ shape functions in their
one-dimensional analogue, where the picture is cleanest. On the reference interval
$[0,1]$ the two $P_1$ hats are the straight lines $1-\hat x$ and $\hat x$; the three
$P_2$ functions are the parabolas $(1-\hat x)(1-2\hat x)$, $\hat x(2\hat x-1)$, and
the bump $4\hat x(1-\hat x)$ on the midpoint. In every case the cardinality property
is visible: each curve passes through $1$ at its own node and $0$ at the others.

\begin{figure}[t]
\centering
\begin{tikzpicture}
% --- P1 panel ---
\begin{axis}[
  name=p1, width=6.4cm, height=4.4cm,
  xlabel={$\hat x$}, ylabel={},
  title={\footnotesize $P_1$ shape functions (1-D)},
  xmin=0, xmax=1, ymin=-0.15, ymax=1.15,
  xtick={0,0.5,1}, ytick={0,0.5,1},
  tick label style={font=\scriptsize}, label style={font=\scriptsize},
  title style={font=\footnotesize}, axis lines=left, every axis plot/.append style={thick},
]
  \addplot[simBlue,  domain=0:1, samples=2] {1-x};
  \addplot[simRust,  domain=0:1, samples=2] {x};
  \node[font=\scriptsize,simBlue] at (axis cs:0.18,0.92) {$\hat\phi_0$};
  \node[font=\scriptsize,simRust] at (axis cs:0.82,0.92) {$\hat\phi_1$};
  \addplot[only marks, mark=*, mark size=1.2pt, black] coordinates {(0,1)(1,0)(0,0)(1,1)};
\end{axis}
% --- P2 panel ---
\begin{axis}[
  name=p2, at={(p1.east)}, anchor=west, xshift=14mm,
  width=6.4cm, height=4.4cm,
  xlabel={$\hat x$}, ylabel={},
  title={\footnotesize $P_2$ shape functions (1-D)},
  xmin=0, xmax=1, ymin=-0.35, ymax=1.15,
  xtick={0,0.5,1}, ytick={0,0.5,1},
  tick label style={font=\scriptsize}, label style={font=\scriptsize},
  title style={font=\footnotesize}, axis lines=left, every axis plot/.append style={thick},
]
  \addplot[simBlue,   domain=0:1, samples=60] {(1-x)*(1-2*x)};
  \addplot[simRust,   domain=0:1, samples=60] {x*(2*x-1)};
  \addplot[simGreen,  domain=0:1, samples=60] {4*x*(1-x)};
  \node[font=\scriptsize,simBlue]  at (axis cs:0.12,0.95) {$\hat\phi_0$};
  \node[font=\scriptsize,simRust]  at (axis cs:0.88,0.95) {$\hat\phi_1$};
  \node[font=\scriptsize,simGreen] at (axis cs:0.5,1.02)  {$\hat\phi_{01}$};
  \addplot[only marks, mark=*, mark size=1.2pt, black] coordinates {(0,1)(0.5,0)(1,0)(0,0)(0.5,1)(1,1)};
\end{axis}
\end{tikzpicture}
\caption[$P_1$ and $P_2$ shape functions]{Lagrange shape functions, shown in their
one-dimensional analogue where the curves are uncluttered. \emph{Left:} the two
$P_1$ ``hat'' functions, $\hat\phi_0=1-\hat x$ and $\hat\phi_1=\hat x$, each linear,
each one at its own node ($\hat x=0$ and $\hat x=1$) and zero at the other.
\emph{Right:} the three $P_2$ functions on $[0,1]$: two vertex parabolas
$\hat\phi_0=(1-\hat x)(1-2\hat x)$, $\hat\phi_1=\hat x(2\hat x-1)$, and the
edge-midpoint bump $\hat\phi_{01}=4\hat x(1-\hat x)$ centred at $\hat x=\tfrac12$.
The dots mark the cardinality property \eqref{eq:cardinality}: every shape function
passes through $1$ at its own node and $0$ at all the others. On the triangle the
same functions appear in barycentric form, $\lambda_i(2\lambda_i-1)$ and
$4\lambda_i\lambda_j$.}
\label{fig:shapefns}
\end{figure}

The pattern continues to any order. Raising $p$ is exactly the $p$-refinement of
\cref{ch:functionspaces}: more shape functions per cell, a richer polynomial space,
faster convergence on smooth fields---and, in Palace, the family of collections
\code{H1\_FECollection}, \code{ND\_FECollection}, and the rest selecting both the de
Rham slot and the order. The point for this chapter is structural: \emph{whatever}
the order, the shape functions live on the reference cell and are tabulated once.

\section{The geometric map and how it transforms calculus}
\label{sec:map}

The reference cell carries the polynomials; the mesh carries the geometry. The bridge
is the geometric map $F_K$, and the heart of this section is what that map does to
the three operations a weak form needs: \emph{evaluating} a field, \emph{differentiating}
it, and \emph{integrating} it. Get the map's bookkeeping right and the same reference
tabulation serves cells of every size, shape, and orientation.

\subsection{The map and its Jacobian}

For each physical cell $K$ the geometric map is a function
\[
  \vx = F_K(\hat\vx)
\]
taking reference coordinates $\hat\vx\in\hat K$ to physical coordinates $\vx\in K$.
For a straight-sided (\emph{affine}) simplex the map is itself built from the nodal
shape functions: writing the physical vertices as $\vx_0,\vx_1,\vx_2$,
\[
  F_K(\hat\vx) = \sum_i \vx_i\,\hat\phi_i(\hat\vx)
  = \vx_0 + (\vx_1-\vx_0)\,\hat x + (\vx_2-\vx_0)\,\hat y .
\]
The map is \emph{affine}: a constant matrix times $\hat\vx$ plus a shift. Its
derivative, the \emph{Jacobian}, is the constant matrix
\[
  J = \frac{\partial F_K}{\partial\hat\vx}
    = \bigl[\;\vx_1-\vx_0 \quad \vx_2-\vx_0\;\bigr]
    \in\Real^{d\times d},
\]
whose columns are the physical edge vectors emanating from vertex $0$. (When the cell
is \emph{curved}---a high-order ``isoparametric'' element whose faces bend to follow
a curved boundary---the map is built from higher-order shape functions and $J$ varies
from point to point. Palace handles both; the affine case is the one to hold in your
head, and the one this chapter computes.) The Jacobian is the single object that
encodes how this particular cell differs from the reference: its determinant measures
how the cell's \emph{volume} compares to the reference, and its inverse-transpose
measures how \emph{directions} are bent.

\subsection{Three transformation rules}

A scalar field is the simplest case. A physical basis function is just the reference
shape function read in reference coordinates:
\[
  \phi_i(\vx) = \hat\phi_i\bigl(F_K^{-1}(\vx)\bigr),
  \qquad\text{equivalently}\qquad
  \phi_i\bigl(F_K(\hat\vx)\bigr) = \hat\phi_i(\hat\vx).
\]
\emph{Values do not change}: to evaluate $\phi_i$ at a physical point, find its
reference preimage and read the table. This is why basis values need to be tabulated
only once.

Derivatives are where the map bites. The weak forms of \cref{ch:weak} integrate
\emph{gradients} of basis functions---the $\Grad\phi_i\cdot\Grad\phi_j$ of the
stiffness matrix---and a gradient taken in physical coordinates is not the gradient
taken in reference coordinates. By the chain rule, differentiating
$\hat\phi_i(\hat\vx)=\phi_i(F_K(\hat\vx))$ with respect to $\hat\vx$,
\[
  \Grad_{\hat\vx}\hat\phi_i = J^{\!\top}\,\Grad_{\vx}\phi_i,
  \qquad\text{hence}\qquad
  \boxed{\;\Grad_{\vx}\phi_i = J^{-\top}\,\Grad_{\hat\vx}\hat\phi_i\;}.
\]
The physical gradient is the \emph{inverse-transpose} Jacobian applied to the
reference gradient. The reference gradient $\Grad_{\hat\vx}\hat\phi_i$ is tabulated
once (for $P_1$ it is the constant $\Grad\lambda_i$ we already wrote down); the cell
contributes only the small matrix $J^{-\top}$. \Cref{fig:jacobian} draws the picture:
the inverse-transpose is precisely the operator that keeps a gradient pointing
``across'' the level sets of $\phi_i$ after the cell is stretched and skewed.

Finally, integration. The change-of-variables theorem turns a physical integral into
a reference integral, paying a \emph{volume factor} $|\det J|$:
\[
  \int_K g(\vx)\,d\vx = \int_{\hat K} g\bigl(F_K(\hat\vx)\bigr)\,\abs{\det J}\,d\hat\vx .
\]
The factor $|\det J|$ is the ratio of the physical cell's volume to the reference
cell's: for the affine triangle, $|\det J| = 2\,\mathrm{area}(K)$ (since the reference
triangle has area $\tfrac12$). Every matrix entry in the book is an integral of this
form, and the recipe is always: \emph{pull back to the reference cell, multiply by
$|\det J|$, and integrate there}---where the basis values are already tabulated.

\begin{figure}[t]
\centering
\begin{tikzpicture}[scale=1.0, every node/.style={font=\scriptsize}]
% reference: gradient perpendicular to level sets
\begin{scope}
  \node[font=\footnotesize\bfseries] at (1.0,2.4) {reference $\hat K$};
  \fill[simBgBlue] (0,0)--(2,0)--(0,2)--cycle;
  \draw[simInk,thick] (0,0)--(2,0)--(0,2)--cycle;
  % level sets of lambda1 = xhat (vertical lines)
  \foreach \x in {0.5,1.0,1.5}{\draw[simGray,densely dotted] (\x,0)--(\x,{2-\x});}
  % gradient arrow (points +x)
  \draw[-{Stealth[length=2.4mm]},simBlue,very thick] (0.7,0.5)--(1.5,0.5);
  \node[simBlue] at (1.1,0.78) {$\Grad_{\hat\vx}\hat\phi$};
  \node[simGray] at (1.45,1.55) {level sets};
\end{scope}
% the J^{-T} map arrow
\draw[flow,simRust,thick] (2.5,1.0) .. controls (3.3,1.4) and (3.9,1.4) .. (4.7,1.0)
  node[midway,above,font=\scriptsize]{$J^{-\top}$};
% physical: skewed level sets, gradient still perpendicular to THEM
\begin{scope}[shift={(5.3,0)}]
  \node[font=\footnotesize\bfseries] at (1.2,2.4) {physical $K$};
  \fill[simBgRust] (0,0)--(2.4,0.4)--(0.7,2.0)--cycle;
  \draw[simInk,thick] (0,0)--(2.4,0.4)--(0.7,2.0)--cycle;
  % skewed level sets
  \draw[simGray,densely dotted] (0.6,0.1)--(0.25,1.1);
  \draw[simGray,densely dotted] (1.2,0.2)--(0.5,1.55);
  \draw[simGray,densely dotted] (1.8,0.3)--(0.78,1.85);
  % gradient: must stay perpendicular to the SKEWED level sets (so it tilts)
  \draw[-{Stealth[length=2.4mm]},simBlue,very thick] (0.8,0.55)--(1.55,0.78);
  \node[simBlue] at (1.25,1.05) {$\Grad_{\vx}\phi$};
  \node[simGray] at (1.6,1.6) {level sets bent};
\end{scope}
\end{tikzpicture}
\caption[The inverse-transpose Jacobian transforms gradients]{Why the gradient
transforms by $J^{-\top}$, not by $J$. \emph{Left:} on the reference cell the level
sets of a shape function $\hat\phi$ (dotted) are evenly spaced and the gradient (blue
arrow) is perpendicular to them. \emph{Right:} the map skews the cell and bends the
level sets; the physical gradient must remain perpendicular to the \emph{bent} level
sets. The operator that does this is the inverse-transpose Jacobian $J^{-\top}$,
applied to the tabulated reference gradient. Tabulate $\Grad_{\hat\vx}\hat\phi$ once;
each cell supplies only its small matrix $J^{-\top}$. In \cref{ch:matrixfree} this
$J^{-\top}$, combined with $|\det J|$ and the material coefficient, is exactly the
per-quadrature-point geometry factor that the $\mathsf{D}$ stage applies.}
\label{fig:jacobian}
\end{figure}

\subsection{The geometry factor: assembling the metric}

Collect the pieces and a clean pattern emerges. The two integrals every solver in
this book needs are the \emph{mass} term, $\int_K \phi_i\,\phi_j\,d\vx$, and the
\emph{stiffness} (grad--grad) term, $\int_K \Grad\phi_i\cdot\Grad\phi_j\,d\vx$.
Pulling each back to the reference cell:
\[
  \int_K \phi_i\phi_j\,d\vx
  = \int_{\hat K} \hat\phi_i\hat\phi_j\;\underbrace{\abs{\det J}}_{\text{mass factor}}\,d\hat\vx,
\]
\[
  \int_K \Grad\phi_i\cdot\Grad\phi_j\,d\vx
  = \int_{\hat K} \bigl(J^{-\top}\Grad_{\hat\vx}\hat\phi_i\bigr)\cdot
                  \bigl(J^{-\top}\Grad_{\hat\vx}\hat\phi_j\bigr)\,\abs{\det J}\,d\hat\vx
  = \int_{\hat K} \Grad_{\hat\vx}\hat\phi_i\cdot
      \underbrace{\bigl(J^{-1}J^{-\top}\abs{\det J}\bigr)}_{\text{metric }G}\,
      \Grad_{\hat\vx}\hat\phi_j\;d\hat\vx .
\]
The mass integral needs only the scalar $|\det J|$; the stiffness integral needs the
$d\times d$ symmetric matrix $G = J^{-1}J^{-\top}\,|\det J|$, the \emph{metric}
introduced by the map. These two objects---$|\det J|$ for mass terms,
$J^{-1}J^{-\top}|\det J|$ for grad--grad terms---are \emph{all} the geometry a cell
contributes. Everything else is reference data.

\begin{keyidea}
A cell contributes to the integrals only through a small \emph{geometry factor}: the
scalar $|\det J|$ for mass terms and the matrix $G = J^{-1}J^{-\top}\,|\det J|$ for
grad--grad terms. The basis values and reference gradients are tabulated once and
shared by every cell; the cell supplies only its Jacobian-derived factor. This is
precisely the \code{geom\_factor\_build} stage of the matrix-free pipeline: it
precomputes $G$ (or $|\det J|$), per quadrature point, once per mesh, and the
material coefficient is pre-multiplied into it so the run-time apply is a single
pointwise multiply.
\end{keyidea}

This factorization is the content of Palace's \code{geom\_factor\_build}: a
\emph{setup-stratum} pass that computes the geometry metric at every quadrature point
once per mesh and stores it, so that the per-apply work---which happens thousands of
times inside a Krylov solve---is reduced to multiplying by a stored number. We meet
that stratification again, in full, in \cref{ch:matrixfree}; the point here is that it
falls straight out of the change-of-variables rules above. On an affine cell $J$ is
constant, so $G$ is constant over the cell and the stored factor is one matrix per
cell; on a curved cell $J$ varies and $G$ is stored per quadrature point.

\section{Quadrature: integrals as weighted sums}
\label{sec:quad}

We can now pull every integral back to the reference cell. One step remains: actually
\emph{evaluating} the reference integral $\int_{\hat K} g(\hat\vx)\,d\hat\vx$. Except
in the simplest cases we do not compute it in closed form. We approximate it by a
\emph{quadrature rule}---a finite weighted sum of the integrand sampled at a few
chosen points.

\subsection{The quadrature idea}

A quadrature rule on $\hat K$ is a set of \emph{points} $\hat\vx_q\in\hat K$ and
\emph{weights} $w_q$ such that
\begin{equation}
  \int_{\hat K} g(\hat\vx)\,d\hat\vx \;\approx\; \sum_{q=1}^{Q} w_q\,g(\hat\vx_q).
  \label{eq:quad}
\end{equation}
The integral---an infinite, continuous operation---is replaced by a finite sum: sample
the integrand at the $Q$ quadrature points, scale each sample by its weight, add them
up. The art is in choosing the points and weights so the sum is accurate. The
classical answer is \emph{Gauss quadrature}: for a given number of points, place them
and weight them so that the rule integrates polynomials \emph{exactly} up to as high
a degree as possible.

\begin{definition}[Exactness degree]
A quadrature rule has \emph{exactness degree} $r$ if \eqref{eq:quad} is an exact
equality for every polynomial $g$ of total degree $\le r$, and fails for some
polynomial of degree $r+1$. A $Q$-point Gauss rule on an interval achieves exactness
degree $2Q-1$: $Q$ points buy $2Q$ free parameters (the points and the weights), and
$2Q$ parameters pin down a polynomial of degree $2Q-1$.
\end{definition}

Exactness degree is the contract a quadrature rule offers, and it is exactly what a
finite-element assembler needs to honor. The integrand of a mass-matrix entry is
$\hat\phi_i\hat\phi_j\,|\det J|$; on $P_p$ elements with an affine map ($|\det J|$
constant) this is a polynomial of degree $2p$. To integrate it without error the rule
must have exactness degree at least $2p$. Choose the rule too weak and the mass matrix
is computed inaccurately---\emph{under-integrated}, a real and sometimes deliberate
choice; choose it stronger than needed and you simply pay for extra points
(\emph{over-integration}, sometimes used to tame variable coefficients). The
exactness degree is the dial.

\begin{figure}[t]
\centering
\begin{tikzpicture}[scale=1.55, every node/.style={font=\scriptsize}]
% three reference triangles with quadrature points
\foreach \dx/\ttl/\desc in {0/{1-point}/{degree 1}, 2.7/{3-point}/{degree 2}, 5.4/{6-point}/{degree 4}}{
  \begin{scope}[shift={(\dx,0)}]
    \fill[simBgBlue] (0,0)--(1.4,0)--(0,1.4)--cycle;
    \draw[simInk,thick] (0,0)--(1.4,0)--(0,1.4)--cycle;
    \node[font=\footnotesize\bfseries] at (0.6,1.62) {\ttl};
    \node[simGray] at (0.6,-0.3) {\desc};
  \end{scope}
}
% 1-point: centroid
\fill[simRust] (0.467,0.467) circle (1.6pt);
% 3-point: at (1/6,1/6),(2/3,1/6),(1/6,2/3) scaled by 1.4
\foreach \px/\py in {0.233/0.233, 0.933/0.233, 0.233/0.933}{
  \fill[simRust] (2.7+\px,\py) circle (1.4pt);}
% 6-point (degree 4): two orbits — approximate symmetric points
\foreach \px/\py in {0.63/0.63, 0.13/0.63, 0.63/0.13, 1.04/0.18, 0.18/1.04, 0.18/0.18}{
  \fill[simRust] (5.4+\px,\py) circle (1.2pt);}
\end{tikzpicture}
\caption[Gauss quadrature points on the reference triangle]{Symmetric Gauss
quadrature rules on the reference triangle, in order of increasing exactness.
\emph{Left:} the $1$-point rule places a single node at the centroid with weight
equal to the triangle's area; it is exact to degree $1$ (it integrates linear
functions exactly). \emph{Middle:} the $3$-point rule (nodes near the edge midpoints,
each weight one-third of the area) is exact to degree $2$. \emph{Right:} a $6$-point
rule, exact to degree $4$, suitable for $P_2$ mass matrices. More points buy higher
exactness; the assembler picks the smallest rule whose exactness degree covers the
integrand's polynomial degree. These points and weights are the discrete realization
of the integral $\int_{\hat K}(\cdot)\,d\hat\vx$, and they are the $P$ axis of the
element-local tensor (\cref{sec:tensor}).}
\label{fig:quadpts}
\end{figure}

\Cref{fig:quadpts} shows three symmetric rules on the reference triangle. The
$1$-point rule samples only the centroid; the $3$-point rule samples three interior
points; a $6$-point rule reaches exactness degree $4$. On the triangle the weights are
fractions of the reference area $\tfrac12$, and the points are placed symmetrically so
that the rule respects the triangle's symmetry. The quadrature points and weights are,
quite literally, the discrete realization of the integral: once the rule is chosen,
``integrate over the cell'' \emph{means} ``sample at these points and take this
weighted sum,'' nothing more.

\begin{physicsbox}
Quadrature is where the continuous integral of the weak form becomes a finite
computation a machine can run. In the matrix-free pipeline of \cref{ch:matrixfree},
the quadrature points are the $P$ axis, the basis tabulation $\mathsf{B}$ evaluates
the trial field \emph{at} those points, the geometry-and-material factor $\mathsf{D}$
multiplies pointwise \emph{at} those points, and the transpose $\mathsf{B}^{\!\top}$
sums the contributions back with the weights folded in. The entire operator apply is
organized around the quadrature points: they are the grain at which the discrete
physics is evaluated.
\end{physicsbox}

\subsection{Tying quadrature to the matrix-free \texttt{D} stage}

It is worth being explicit about where quadrature lands in the pipeline, because it
is the conceptual hinge between this chapter and \cref{ch:matrixfree}. An operator
apply $\vy = \mathsf{A}\vx$ in the matrix-free scheme is the composition
\[
  \mathsf{A} = \mathsf{G}^{\!\top}\,\mathsf{B}^{\!\top}\,\mathsf{D}\,\mathsf{B}\,\mathsf{G},
\]
read right to left: $\mathsf{G}$ gathers the global coefficient vector into
per-element local tensors; $\mathsf{B}$ applies the tabulated reference basis to get
field values at the quadrature points; $\mathsf{D}$ multiplies, point by point, by the
geometry-and-material factor (the $|\det J|$ or $G$ of \cref{sec:map}, with the
quadrature weight $w_q$ and the material coefficient pre-multiplied in); $\mathsf{B}^{\!\top}$
contracts the weighted point values back against the basis; $\mathsf{G}^{\!\top}$
scatters to the global vector. The weighted sum \eqref{eq:quad} is split across two
stages: the weight $w_q$ rides inside $\mathsf{D}$, and the summation $\sum_q$ is the
contraction that $\mathsf{B}^{\!\top}$ performs. Quadrature is not a separate step
bolted on---it \emph{is} the structure of the apply.

\section{Vector elements and the Piola maps}
\label{sec:vector}

The nodal elements of \cref{sec:shape} are the scalar case. The vector spaces
$\Hcurl$ and $\Hdiv$ of \cref{ch:functionspaces}---the homes of $\vE$, $\vA$, and the
fluxes $\vB$, $\vD$---use \emph{vector-valued} shape functions, and they map to the
physical cell by a different rule. We flag the essentials here and leave the full
edge-basis machinery to the references; the point is that the reference-element frame
carries over intact, with one twist in the transformation.

A N\'ed\'elec ($\Hcurl$) reference shape function is a vector field on $\hat K$---the
Whitney function $\hat\phi_{01}=\lambda_0\Grad\lambda_1-\lambda_1\Grad\lambda_0$ of
\cref{wex:whitney} is the lowest-order example. A Raviart--Thomas ($\Hdiv$) reference
function is likewise vector-valued. The subtlety is that mapping a vector field to a
skewed physical cell by the \emph{value-preserving} rule of the scalar case would
\emph{destroy} the tangential or normal continuity that \cref{ch:functionspaces}
showed is the whole point of these spaces. A scalar value is the same in any
coordinates; a vector's tangential and normal parts are not.

The fix is the \emph{Piola transforms}, which are designed to preserve exactly the
right component:
\[
  \text{covariant (}\Hcurl\text{):}\quad \vu(\vx) = J^{-\top}\hat\vu(\hat\vx),
  \qquad
  \text{contravariant (}\Hdiv\text{):}\quad \vu(\vx) = \frac{1}{\det J}\,J\,\hat\vu(\hat\vx).
\]
The covariant (curl-conforming) map applies $J^{-\top}$---the same operator that
transforms gradients---and it preserves \emph{tangential} components along edges; this
is why edge circulations survive the map and tangential continuity is maintained
cell to cell. The contravariant (divergence-conforming) map applies $J/\det J$ and
preserves \emph{normal} components through faces, maintaining the flux continuity of
$\Hdiv$.

\begin{keyidea}
The reference-element idea extends to the vector spaces, but with a
\emph{component-preserving} map in place of the value-preserving one. The covariant
Piola map $J^{-\top}$ preserves tangential traces (so $\Hcurl$ edge circulations
survive) and the contravariant Piola map $J/\det J$ preserves normal traces (so
$\Hdiv$ face fluxes survive). The continuity that \cref{ch:functionspaces} built into
the DoFs is preserved by the map---which is what makes a single edge or face DoF,
shared by two cells, mean the same physical quantity on both.
\end{keyidea}

The remarkable fact---which \cref{ch:derham}'s de Rham complex explains---is that the
three maps fit together: nodal values transform trivially, gradients and $\Hcurl$
fields transform by $J^{-\top}$, $\Hdiv$ fields by $J/\det J$, and the chain
$\Grad\to\Curl\to\Div$ commutes with the maps. For this chapter the operative point is
narrow and practical: a vector element is still a reference shape function tabulated
once, carried to each cell by a Jacobian-built map. Only the map's formula changes.

\section{The element-local tensor: $E$, $L$, $P$, $C$, $G$}
\label{sec:tensor}

We close by collecting everything this chapter has produced into the single data
structure the matrix-free pipeline of \cref{ch:matrixfree} runs on. Every quantity
we have met---the local degrees of freedom, the quadrature points, the field
components, the geometry factor---is an axis of a small per-element tensor. Naming
those axes now makes \cref{ch:matrixfree} a matter of contraction bookkeeping rather
than fresh concepts.

Restrict the global field to one element and you have a small vector: one coefficient
for each local DoF of that element. Stacked over all elements, this is a tensor with
two axes,
\[
  \texttt{Tensor[(E, L)]},
\]
where $E$ counts the mesh elements and $L$ counts the local DoFs per element (for
$P_1$ on a triangle, $L=3$; the reference shape functions of \cref{sec:shape}).
Apply the tabulated basis $\mathsf{B}$---evaluate those $L$ shape functions at the
$P$ quadrature points---and you get the field's values at every quadrature point of
every element,
\[
  \texttt{Tensor[(E, P, C)]},
\]
where $P$ counts quadrature points per element (the rule of \cref{sec:quad}) and $C$
counts value components ($C=1$ for a scalar field; $C=d$ for a gradient or a vector
field). Finally the geometry factor of \cref{sec:map}---the per-point $|\det J|$ or
metric $G$, with weight and material pre-multiplied in---is itself a per-point tensor,
\[
  \texttt{Tensor[(E, P, G)]},
\]
with $G$ counting the stored geometry-metric components. These three tensors, over the
five named axes $E,L,P,C,G$, are the complete vocabulary of element-local finite
element computation.

\begin{figure}[t]
\centering
\begin{tikzpicture}[scale=1.0, every node/.style={font=\scriptsize}]
% [E,L] block
\node[draw, simBlue, fill=simBgBlue, thick, minimum width=20mm, minimum height=13mm,
  align=center] (el) at (0,0) {\code{Tensor}\\\code{[(E, L)]}};
\node[font=\scriptsize, simGray, align=center] at (0,-1.35)
  {per-element\\local DoFs};
% B arrow
\draw[flow, simInk] (1.1,0)--(3.0,0)
  node[midway, above, font=\footnotesize] {$\mathsf{B}$}
  node[midway, below, font=\scriptsize, simGray] {tabulated basis};
% [E,P,C] block
\node[draw, simRust, fill=simBgRust, thick, minimum width=22mm, minimum height=13mm,
  align=center] (epc) at (4.2,0) {\code{Tensor}\\\code{[(E, P, C)]}};
\node[font=\scriptsize, simGray, align=center] at (4.2,-1.35)
  {values at\\quad points};
% D arrow (with geom factor coming in)
\draw[flow, simInk] (5.4,0)--(7.3,0)
  node[midway, above, font=\footnotesize] {$\mathsf{D}$}
  node[midway, below, font=\scriptsize, simGray] {pointwise};
% geom factor feeding D
\node[draw, simGreen, fill=simBgGreen, semithick, minimum width=20mm, minimum height=9mm,
  align=center] (epg) at (6.35,1.5) {\code{Tensor[(E, P, G)]}};
\node[font=\scriptsize, simGreen] at (6.35,2.15) {geometry factor};
\draw[depedge, simGreen] (epg.south)--(6.35,0.42);
% output
\node[draw, simRust, fill=simBgRust, thick, minimum width=22mm, minimum height=13mm,
  align=center] (out) at (8.5,0) {\code{Tensor}\\\code{[(E, P, C)]}};
\node[font=\scriptsize, simGray, align=center] at (8.5,-1.35)
  {weighted\\values};
\end{tikzpicture}
\caption[The element-local tensor and its axes]{The element-local tensor pipeline,
the data structure \cref{ch:matrixfree} runs on. A field restricted to elements is a
\code{Tensor[(E, L)]} of per-element local DoFs ($E$ elements, $L$ local DoFs). The
tabulated basis $\mathsf{B}$ of \cref{sec:shape}---the reference shape functions
evaluated at the quadrature points---contracts it to a \code{Tensor[(E, P, C)]} of
field values at the $P$ quadrature points ($C$ components). The pointwise stage
$\mathsf{D}$ multiplies by the geometry-and-material factor
\code{Tensor[(E, P, G)]}---the $|\det J|$ or metric $G$ of \cref{sec:map}, with
quadrature weight and material coefficient pre-multiplied ($G$ metric components).
The five axes $E$, $L$, $P$, $C$, $G$ are exactly the quantities this chapter built:
elements, reference DoFs, quadrature points, field components, and geometry-metric
storage.}
\label{fig:elemtensor}
\end{figure}

\Cref{fig:elemtensor} draws the flow. Read it against the three sections that built
the pieces: the $L$ axis is the reference shape functions of \cref{sec:shape}; the
$P$ axis is the quadrature rule of \cref{sec:quad}; the $C$ axis is the value or
derivative components selected by the weak-form differential operator; the $G$ axis is
the geometry factor of \cref{sec:map}; and $E$ simply ranges over the cells of the
mesh, each carrying its own Jacobian but sharing the one reference tabulation. This is
the payoff of the reference-element idea stated as a data structure: the expensive,
shared content (the basis) is the small $\mathsf{B}$ table indexed by $(L,P)$, and the
cell-specific content (the geometry) is the $(E,P,G)$ factor---built once, by
\code{geom\_factor\_build}, and reused across every apply.

\begin{notationbox}
The five element-local axes recur throughout the framework chapters:
$E$ (mesh elements), $L$ (local DoFs per element), $P$ (quadrature points per
element), $C$ (field/derivative components), $G$ (geometry-metric components). The
global true-DoF axis $N$ of \cref{ch:functionspaces} is \emph{not} one of
them---$N$ is the flat global vector, and the element-restriction operator
$\mathsf{G}$ is precisely the boundary between the global $N$ vocabulary and the
element-local $(E,L)$ vocabulary. \Cref{ch:matrixfree} works entirely in these axes.
\end{notationbox}

\section{Worked examples}

The two examples carry the chapter's arithmetic to the end. The first integrates a
product of shape functions by quadrature and checks it against the exact mass-matrix
entry, exhibiting the exactness-degree contract concretely. The second computes a
Jacobian and transforms a gradient, exhibiting the map's three rules on a specific
physical triangle.

\subsection{A mass-matrix entry by quadrature}

\begin{workedexample}{Integrating two $P_1$ hats: exact, 1-point, and 3-point}{massmatrix}
Work on the reference triangle $\hat K$ with vertices $(0,0),(1,0),(0,1)$ and area
$\tfrac12$. The three $P_1$ shape functions are the barycentric coordinates
\[
  \hat\phi_0 = 1-\hat x-\hat y, \qquad \hat\phi_1 = \hat x, \qquad \hat\phi_2 = \hat y .
\]
We compute the reference mass-matrix entry
$m_{01} = \int_{\hat K}\hat\phi_0\hat\phi_1\,d\hat\vx
        = \int_{\hat K}(1-\hat x-\hat y)\,\hat x\,d\hat\vx$
three ways.

\smallskip\noindent\textbf{Exact.} There is a standard formula for integrating
products of barycentric coordinates over a triangle of area $A$:
\[
  \int_{\hat K}\lambda_0^{a}\lambda_1^{b}\lambda_2^{c}\,d\hat\vx
  = 2A\,\frac{a!\,b!\,c!}{(a+b+c+2)!}.
\]
With $A=\tfrac12$ and $(a,b,c)=(1,1,0)$ this gives
$m_{01} = 2\cdot\tfrac12\cdot\dfrac{1!\,1!\,0!}{4!} = \dfrac{1}{24}\approx 0.041667$.
(For comparison, the diagonal entry with $(a,b,c)=(2,0,0)$ is
$2\cdot\tfrac12\cdot\tfrac{2!}{4!}=\tfrac{1}{12}$.) So the exact answer is
$m_{01} = \tfrac{1}{24}$.

\smallskip\noindent\textbf{1-point rule (centroid, degree 1).} The one-point Gauss
rule samples the centroid $\hat\vx_c=(\tfrac13,\tfrac13)$ with weight equal to the
area, $w=\tfrac12$. The integrand there is
$\hat\phi_0\hat\phi_1 = (1-\tfrac13-\tfrac13)\cdot\tfrac13 = \tfrac13\cdot\tfrac13 = \tfrac19$,
so
\[
  m_{01}^{(1)} = w\,\hat\phi_0(\hat\vx_c)\hat\phi_1(\hat\vx_c)
              = \tfrac12\cdot\tfrac19 = \tfrac{1}{18} \approx 0.055556 .
\]
This is \emph{wrong}---off by a third. The reason is the exactness contract: the
integrand $\hat\phi_0\hat\phi_1$ is a quadratic (degree $2$), but the $1$-point rule
is exact only to degree $1$. Under-integration shows up as a concrete numerical error.

\smallskip\noindent\textbf{3-point rule (degree 2).} The symmetric three-point rule
places nodes at $(\tfrac16,\tfrac16)$, $(\tfrac23,\tfrac16)$, $(\tfrac16,\tfrac23)$,
each with weight $w_q=\tfrac16$ (one-third of the area $\tfrac12$). Evaluate the
integrand $\hat\phi_0\hat\phi_1=(1-\hat x-\hat y)\,\hat x$ at each:
\begin{align*}
  q=1:\ (\tfrac16,\tfrac16) &: (1-\tfrac13)\cdot\tfrac16 = \tfrac23\cdot\tfrac16 = \tfrac{1}{9}, \\
  q=2:\ (\tfrac23,\tfrac16) &: (1-\tfrac56)\cdot\tfrac23 = \tfrac16\cdot\tfrac23 = \tfrac{1}{9}, \\
  q=3:\ (\tfrac16,\tfrac23) &: (1-\tfrac56)\cdot\tfrac16 = \tfrac16\cdot\tfrac16 = \tfrac{1}{36}.
\end{align*}
The weighted sum is
\[
  m_{01}^{(3)} = \tfrac16\Bigl(\tfrac19+\tfrac19+\tfrac1{36}\Bigr)
              = \tfrac16\cdot\Bigl(\tfrac{4+4+1}{36}\Bigr)
              = \tfrac16\cdot\tfrac{9}{36} = \tfrac16\cdot\tfrac14 = \tfrac{1}{24}.
\]
This is \emph{exact}. The three-point rule has exactness degree $2$, the integrand
has degree $2$, and the contract is honored: the quadrature equals the closed-form
integral to the last digit. The lesson is the practical rule of \cref{sec:quad}: to
integrate a $P_p$ mass matrix exactly on an affine cell, use a rule of exactness
degree $\ge 2p$. Here $p=1$, $2p=2$, and the three-point (degree-$2$) rule is the
smallest that suffices---the one-point rule, degree $1<2$, is not enough.
\end{workedexample}

\subsection{A Jacobian and a transformed gradient}

\begin{workedexample}{The map of one physical triangle, and a gradient through it}{jacobian}
Map the reference triangle $\hat K$ (vertices $(0,0),(1,0),(0,1)$) onto the physical
triangle $K$ with vertices
\[
  \vx_0 = (1,1), \qquad \vx_1 = (4,2), \qquad \vx_2 = (2,5).
\]
\textbf{The Jacobian.} The affine map is
$F_K(\hat\vx) = \vx_0 + (\vx_1-\vx_0)\hat x + (\vx_2-\vx_0)\hat y$, so the Jacobian is
the constant matrix whose columns are the physical edge vectors from $\vx_0$:
\[
  J = \bigl[\,\vx_1-\vx_0 \ \ \vx_2-\vx_0\,\bigr]
    = \begin{bmatrix} 3 & 1 \\ 1 & 4 \end{bmatrix}.
\]
Its determinant is $\det J = 3\cdot 4 - 1\cdot 1 = 11$, so $|\det J| = 11$, and indeed
the area of $K$ is $\tfrac12|\det J| = \tfrac{11}{2}$, eleven times the reference
area $\tfrac12$. The volume factor for any mass integral on this cell is $11$.

\smallskip\noindent\textbf{The inverse-transpose.} The inverse Jacobian is
\[
  J^{-1} = \frac{1}{11}\begin{bmatrix} 4 & -1 \\ -1 & 3 \end{bmatrix},
  \qquad
  J^{-\top} = \frac{1}{11}\begin{bmatrix} 4 & -1 \\ -1 & 3 \end{bmatrix},
\]
which here equals $J^{-1}$ because $J$ is symmetric (an accident of this example;
in general $J^{-\top}\ne J^{-1}$).

\smallskip\noindent\textbf{A transformed gradient.} On the reference cell the shape
function $\hat\phi_1=\hat x$ has constant reference gradient
$\Grad_{\hat\vx}\hat\phi_1 = (1,0)^{\!\top}$. Its physical gradient is
\[
  \Grad_{\vx}\phi_1 = J^{-\top}\Grad_{\hat\vx}\hat\phi_1
  = \frac{1}{11}\begin{bmatrix} 4 & -1 \\ -1 & 3 \end{bmatrix}
    \begin{bmatrix} 1 \\ 0 \end{bmatrix}
  = \frac{1}{11}\begin{bmatrix} 4 \\ -1 \end{bmatrix}
  = \bigl(\tfrac{4}{11},\,-\tfrac{1}{11}\bigr).
\]
Notice it has \emph{tilted}: the reference gradient pointed purely along $\hat x$, but
on the skewed physical cell the gradient of $\phi_1$ has acquired a $y$-component,
exactly as \cref{fig:jacobian} promised---it stays perpendicular to the (now bent)
level sets of $\phi_1$. Likewise $\hat\phi_2=\hat y$ has reference gradient
$(0,1)^{\!\top}$ and physical gradient
$J^{-\top}(0,1)^{\!\top}=(-\tfrac{1}{11},\tfrac{3}{11})$, and
$\hat\phi_0=1-\hat x-\hat y$ has reference gradient $(-1,-1)^{\!\top}$ and physical
gradient $(-\tfrac{3}{11},-\tfrac{2}{11})$. As a check, the three physical gradients
sum to zero---because $\phi_0+\phi_1+\phi_2\equiv 1$ on the cell, a constant, whose
gradient vanishes.

\smallskip\noindent\textbf{Assembling a stiffness entry.} With the gradients in hand,
a grad--grad (stiffness) entry on this affine cell---where $J$, hence everything, is
constant---is just the constant integrand times the cell volume:
\[
  k_{12} = \int_K \Grad\phi_1\cdot\Grad\phi_2\,d\vx
  = \bigl(\Grad_{\vx}\phi_1\cdot\Grad_{\vx}\phi_2\bigr)\,\mathrm{area}(K)
  = \Bigl(\tfrac{4}{11}\cdot(-\tfrac{1}{11}) + (-\tfrac{1}{11})\cdot\tfrac{3}{11}\Bigr)\cdot\tfrac{11}{2}.
\]
The dot product is $\tfrac{-4-3}{121}=-\tfrac{7}{121}$, so
$k_{12} = -\tfrac{7}{121}\cdot\tfrac{11}{2} = -\tfrac{7}{22}\approx -0.318$. No
quadrature was needed: on an affine cell the grad--grad integrand of $P_1$ elements
is constant, so ``integrate'' is just ``multiply by the volume''---the degenerate,
one-effective-point case of the quadrature of \cref{sec:quad}. This single number is
one entry of the element stiffness matrix that \cref{ch:assembly} scatters into the
global operator.
\end{workedexample}

\summarysection
A finite-element code does not build basis functions cell by cell. It defines a small
family of \emph{shape functions} once, on a canonical \emph{reference element}
$\hat K$, and reaches every physical cell $K$ by a geometric map $\vx=F_K(\hat\vx)$.
On the reference simplex the nodal ($\Hone$) shape functions are the barycentric
coordinates $\lambda_i$ (the $P_1$ hats) and their higher-order $P_p$ relatives, each
fixed by the cardinality property $\hat\phi_i(\hat\vx_j)=\delta_{ij}$ that makes a
coefficient equal a nodal value. The map's Jacobian $J$ carries calculus across:
values are unchanged, gradients transform by the inverse-transpose
$\Grad_\vx=J^{-\top}\Grad_{\hat\vx}$, and integrals pay a volume factor $|\det J|$.
The two integrals every solver needs collapse to a per-cell geometry factor---the
scalar $|\det J|$ for mass terms, the metric $G=J^{-1}J^{-\top}|\det J|$ for
grad--grad terms---which is exactly Palace's \code{geom\_factor\_build}. The remaining
integral over $\hat K$ is evaluated by \emph{quadrature}, a weighted sum
$\int_{\hat K}g\approx\sum_q w_q g(\hat\vx_q)$ whose \emph{exactness degree} must
cover the integrand's polynomial degree ($\ge 2p$ for a $P_p$ mass matrix), as the
worked mass-matrix example showed concretely: the one-point rule under-integrates a
quadratic, the three-point rule nails it. The vector spaces $\Hcurl$ and $\Hdiv$ use
the same reference frame with \emph{Piola} maps ($J^{-\top}$ covariant,
$J/\det J$ contravariant) that preserve the tangential and normal continuity built
into their DoFs. Everything collects into the element-local tensor over five axes:
\code{[(E,L)]} per-element DoFs, \code{[(E,P,C)]} values at quadrature points, and
\code{[(E,P,G)]} the geometry factor---the data structure the matrix-free pipeline of
\cref{ch:matrixfree} runs on, with $\mathsf{B}$ the tabulated basis and $\mathsf{D}$
the per-point geometry-and-material multiply.

\furtherreading
The reference-element construction, the Lagrange and higher-order shape functions, the
Jacobian transformation rules, and the quadrature theory of this chapter are developed
rigorously in Brenner and Scott~\citep{brennerscott2008}, the standard graduate text on
the mathematics of finite elements; their treatment of the affine-equivalence of
elements is the precise version of \cref{fig:refmap}. For the engineering view---tables
of shape functions and quadrature rules on triangles, tetrahedra, and hexahedra, with
the electromagnetics applications of this book---Jin~\citep{jin2014} is the natural
companion, and it carries the N\'ed\'elec and Raviart--Thomas vector elements of
\cref{sec:vector} with their Piola maps in full. The matrix-free realization---the
$\mathsf{B}$/$\mathsf{D}$ tensor-contraction pipeline of \cref{sec:tensor}, sum
factorization on tensor-product elements, and the element-local tensor axes---is the
subject of the libCEED project and its reference~\citep{ceed2021}; \cref{ch:matrixfree}
takes it up in detail. The Piola transforms and the way the three maps commute with
$\Grad,\Curl,\Div$ are the bridge to the de Rham complex of \cref{ch:derham}.

\exercisessection
\begin{enumerate}[nosep]
\item Write out the six $P_2$ shape functions on the reference triangle in
  barycentric form ($\lambda_i(2\lambda_i-1)$ for the three vertices and
  $4\lambda_i\lambda_j$ for the three edge midpoints) and verify the cardinality
  property \eqref{eq:cardinality} for one vertex function and one edge function by
  evaluating it at all six nodes (the three vertices and the three edge midpoints).
\item How many shape functions does the $P_3$ space on a triangle have? On a
  tetrahedron the count is $\binom{p+3}{3}$; evaluate it for $p=1,2,3$ and identify
  which DoFs (vertex, edge, face, interior) the new ones at each order sit on.
\item For the reference mass matrix of the $P_1$ triangle, use the barycentric
  integration formula of \cref{wex:massmatrix} to compute every entry $m_{ij}$, and
  show the result is $\tfrac{1}{24}\begin{psmallmatrix}2&1&1\\1&2&1\\1&1&2\end{psmallmatrix}$
  (diagonal $\tfrac{1}{12}$, off-diagonal $\tfrac{1}{24}$). Confirm the row sums equal
  the integrals $\int_{\hat K}\hat\phi_i\,d\hat\vx=\tfrac16$.
\item Repeat the three-way quadrature comparison of \cref{wex:massmatrix} for the
  \emph{diagonal} entry $m_{00}=\int_{\hat K}\hat\phi_0^2\,d\hat\vx$. Compute the exact
  value, the one-point estimate, and the three-point estimate. Does the one-point rule
  do better or worse on the diagonal than on the off-diagonal? Why is the three-point
  rule still exact?
\item Map the reference triangle onto the physical triangle with vertices
  $\vx_0=(0,0)$, $\vx_1=(2,0)$, $\vx_2=(1,3)$. Compute $J$, $\det J$, the area of $K$,
  and the physical gradient $\Grad_\vx\phi_0$ of the shape function
  $\hat\phi_0=1-\hat x-\hat y$. Verify that $\Grad_\vx\phi_0+\Grad_\vx\phi_1+\Grad_\vx\phi_2=\mathbf 0$.
\item Explain why the volume factor for a mass integral is $|\det J|$ but the gradient
  transform involves $J^{-\top}$. Using \cref{wex:jacobian}'s $J$, show by direct
  substitution that $\int_K\Grad\phi_i\cdot\Grad\phi_j\,d\vx$ pulls back to
  $\int_{\hat K}\Grad_{\hat\vx}\hat\phi_i\cdot(J^{-1}J^{-\top}|\det J|)\Grad_{\hat\vx}\hat\phi_j\,d\hat\vx$,
  and identify the metric matrix $G=J^{-1}J^{-\top}|\det J|$ for that cell.
\item A $Q$-point Gauss rule on an interval is exact to degree $2Q-1$. What is the
  minimum number of one-dimensional Gauss points needed to integrate a $P_4$ mass
  matrix exactly (integrand degree $2p=8$) on an affine element? What if the map is
  curved, so that $|\det J|$ is itself a degree-$2$ polynomial---what integrand degree
  must the rule now cover?
\item Identify, for each of the five element-local axes $E,L,P,C,G$ of
  \cref{sec:tensor}, which earlier chapter or section fixes its extent: which is set
  by the mesh, which by the FE order, which by the quadrature rule, which by the
  weak-form differential operator, and which by the space dimension. Then state which
  axes are \emph{build-stratum} (fixed once per mesh and order) and which is
  \emph{run-stratum} (a per-apply field value), connecting to the
  \code{geom\_factor\_build} setup pass.
\end{enumerate}
